\scanpage{87}

% top, heading and first formula
graphe produit de graphes $(\Gamma_i)\ (i\in I)$

% top, vertex set
\[
\text{ens des sommets}\qquad\prod\Gamma_i
\]

% top, adjacency and addition
$x=(x_i)$ et $y=(y_i)$ \deleted{distincts} liés ssi $\exists\,i\in I$ tel que $x_j=y_j$ si $j\ne i$, et $x_i,y_i$ liés
\inserted{c'est op. associative (pas produit catégorique.)}

% upper, opposite graph
graphe \deleted{\unsure{opposé}} $\Gamma^\circ$ \deleted{\unsure{dual}} opposé du graphe $\Gamma$ : $=$ sommets, deux pts \inserted{distincts} de $\Gamma^\circ$ sont liés ssi ils sont non liés dans $\Gamma$.

% upper middle, NB
NB pour un $\Gamma$ donné, l'ens des structures de \deleted{\unsure{la}} graphes \inserted{(strict)} \deleted{\unsure{correspond en bijection par dualité avec l'ens des structures de}} \inserted{$=$} l'ens des parties de \inserted{$^{\wedge}$} $\mathfrak P_2(\Gamma)$, \unsure{la|le} \deleted{\unsure{dualité}} \inserted{opposition} est le passage au complémentaire. C'est une involution qui, pour $\mathfrak P_2(\Gamma)\ne\varnothing$ i.e. $\mathrm{card}(\Gamma)\ge 2$, est sans pt fixe. Néanmoins, il peut arriver que $\Gamma$ soit isomorphe au graphe \deleted{dual}. On dit alors que $\Gamma$ est ``\uline{auto\deleted{dual}\inserted{opposé}}''.

% middle, coproduct definition
\uline{coproduit} de graphes $\Gamma_i$ : \inserted{graphe $\Gamma=\overset{\circ}{\prod}\Gamma_i$} défini par $\Gamma^\circ=\prod_i\Gamma_i^\circ$ \inserted{(notation $\Gamma_1\mathbin{\overset{\circ}{\times}}\Gamma_2$ pour deux facteurs,}
\deleted{$\unsure{\Gamma=}$}
\[
\Gamma=\left(\prod_i\Gamma_i^\circ\right)^\circ\quad\text{i.e.}
\]

% middle-lower, vertices and condition
Donc ens des sommets de $\Gamma$ \deleted{\unsure{est}} $=\prod_i\Gamma_i$. Deux sommets $x=(x_i)$ et $y=(y_i)$ \inserted{de $\Gamma$} sont \inserted{\uline{non}} liés (dans $\Gamma$) ssi ils sont \deleted{non} liés dans $\Gamma^\circ$, i.e. ssi $\exists\,i\in I$ tel que $x_j=y_j$ pour $j\ne i$, et $x_i,y_i$ \deleted{\unsure{non}} liés dans $\Gamma_i^\circ$, i.e. $x_i,y_i$ \uline{non} liés dans $\Gamma_i$. On a prouvé la

% left of condition
\margin{\uline{NB} $\left(\prod\Gamma_i\right)^\circ=\overset{\circ}{\prod}\Gamma_i^\circ$ \deleted{\unsure{donc}}.}

% lower, proposition
\uline{Prop} Un pavé (i.e. \inserted{graphe iso. au} produit de deux triangles) est autodual, en d'autres termes un pavé est isomorphe au coproduit de deux trièdres (duals des triangles) ou encore un coproduit de deux trièdres est un pavé.

% archive sheet number: 86
